\part{Nebenläufigkeit}

\chapter{Goroutinen}\label{kap:goroutinen}

\begin{merksatz}
\boxtitel{Was sie kosten}
\begin{center}
\renewcommand{\arraystretch}{1.5}
\begin{tabular}{@{}P{3.0cm}P{2.2cm}P{2.2cm}P{1.6cm}@{}}
\toprule
& \textbf{System\-faden} & \textbf{Goroutine} & \textbf{Faktor}\\
\midrule
Keller anfangs & $\approx8$ MB & $\mathbf{2}$ \textbf{KB} & $4000$\\
Erzeugen & $\approx12$ µs & $\mathbf{0{,}3}$ µs & $40$\\
Umschalten & $\approx1{,}5$ µs & $\mathbf{0{,}1}$ µs & $15$\\
Praktisch möglich & $\approx10^4$ & $\mathbf{10^6}$ & \\
\bottomrule
\end{tabular}
\end{center}

\begin{center}
\begin{minipage}{0.88\textwidth}
\kursiv{Die erste Zeile ist der Kern: Ein Systemfaden reserviert acht Megabyte Kellerbereich, eine Goroutine zwei Kilobyte, die bei Bedarf wachsen. Deshalb sind eine Million Goroutinen möglich und eine Million Fäden nicht.}
\end{minipage}
\end{center}

\textbf{Was das für die Übertragungsschicht bedeutet.}

\begin{center}
\renewcommand{\arraystretch}{1.5}
\begin{tabular}{@{}P{2.8cm}P{2.6cm}P{2.8cm}@{}}
\toprule
\textbf{Bauform} & \textbf{Bei $1000$ Geräten} & \\
\midrule
Ein Faden je Gerät & \textcolor{akzent}{$\approx8$ GB Keller} & \textcolor{akzent}{unmöglich}\\
Ereignisschleife & \textcolor{gruen}{ein Faden} & \textcolor{akzent}{\textbf{Ablauf zerrissen}}\\
\textbf{Eine Goroutine je Gerät} & \textcolor{gruen}{$\approx\mathbf{20}$ \textbf{MB}} & \textcolor{gruen}{\textbf{Ablauf zusammenhängend}}\\
\bottomrule
\end{tabular}
\end{center}

\begin{center}
\kursiv{Die zweite Zeile ist die Bauform, die man in anderen Sprachen wählen muss: eine Ereignisschleife mit Rückrufen, in der ein logischer Ablauf über zehn Funktionen verteilt liegt. Die dritte Zeile schreibt denselben Ablauf hintereinander --- lesbar wie ein sequentielles Programm.}
\end{center}

\begin{technik}
\boxtitel{Nebenläufigkeit ist nicht Parallelität}
\begin{center}
\renewcommand{\arraystretch}{1.5}
\begin{tabular}{@{}P{2.4cm}P{2.8cm}P{3.2cm}@{}}
\toprule
& \textbf{Nebenläufigkeit} & \textbf{Parallelität}\\
\midrule
Ist & \textbf{eine Entwurfsform} & \textbf{eine Ausführungsform}\\
Betrifft & \textbf{Zerlegung} & Hardware\\
Nützt bei & Warten & \textbf{Rechnen}\\
Auf einem Kern & \textcolor{gruen}{sinnvoll} & \textcolor{akzent}{unmöglich}\\
\bottomrule
\end{tabular}
\end{center}

\begin{center}
\begin{minipage}{0.86\textwidth}
\kursiv{Pikes Unterscheidung ist die nützlichste des Kapitels: Nebenläufigkeit heißt, ein Programm in unabhängig fortschreitende Teile zu zerlegen. Ob diese Teile gleichzeitig laufen, ist eine andere Frage --- und für die Übertragungsschicht, die überwiegend wartet, die unwichtigere.}
\end{minipage}
\end{center}
\end{technik}
\end{merksatz}

\begin{merksatz}
\boxtitel{Wer beendet die Goroutine?}
\begin{center}
\renewcommand{\arraystretch}{1.55}
\begin{tabular}{@{}P{2.8cm}P{6.2cm}@{}}
\toprule
\textbf{Die Regel} & Wer eine Goroutine startet, muss wissen, wann sie endet\\
\textbf{Der Grund} & \textbf{Es gibt keinen Weg, sie von außen abzubrechen}\\
\textbf{Das Mittel} & Ein Abbruchsignal, das sie selbst beachtet\\
\bottomrule
\end{tabular}
\end{center}

\begin{center}
\begin{minipage}{0.88\textwidth}
\kursiv{Die zweite Zeile überrascht viele: Go bietet kein \code{kill} für Goroutinen --- absichtlich, weil ein Abbruch von außen den Zustand in unbekanntem Zustand hinterlässt. Eine Goroutine endet, wenn ihre Funktion endet, und sonst nie.}
\end{minipage}
\end{center}

\begin{center}
\begin{minipage}{0.9\textwidth}
\code{func lauscher(ctx context.Context, g Geraet, aus chan<- Messwert) \{}\\
\code{\ \ \ \ for \{}\\
\code{\ \ \ \ \ \ \ \ m, err := g.Lies()}\\
\code{\ \ \ \ \ \ \ \ if err != nil \{ return \}}\\
\code{\ \ \ \ \ \ \ \ select \{}\\
\code{\ \ \ \ \ \ \ \ case aus <- m:}\\
\code{\ \ \ \ \ \ \ \ case <-ctx.Done():\ \ \ \ // Abbruch beachtet}\\
\code{\ \ \ \ \ \ \ \ \ \ \ \ return}\\
\code{\ \ \ \ \ \ \ \ \}}\\
\code{\ \ \ \ \}}\\
\code{\}}
\end{minipage}
\end{center}

\begin{center}
\kursiv{Der Abbruchkontext ist das erste Argument --- eine Vereinbarung, die durchgehalten wird. Ohne den \code{select} würde die Goroutine ewig auf einen Abnehmer warten, auch wenn niemand mehr zuhört. Das ist das Leck aus Kapitel~\ref{kap:schiefgeht}.}
\end{center}

\begin{warnung}
\boxtitelr{Die drei Enden}
\begin{center}
\renewcommand{\arraystretch}{1.5}
\begin{tabular}{@{}P{2.8cm}P{2.4cm}P{2.8cm}@{}}
\toprule
\textbf{Ende} & \textbf{Sauber?} & \\
\midrule
Funktion endet & \textcolor{gruen}{\textbf{ja}} & \textbf{der Normalfall}\\
Abbruch über Kontext & \textcolor{gruen}{ja} & \\
\code{panic} & \textcolor{akzent}{\textbf{reißt alles mit}} & \textcolor{akzent}{\textbf{das ganze Programm}}\\
Programm endet & \textcolor{mittel}{abrupt} & kein Aufräumen\\
\bottomrule
\end{tabular}
\end{center}

\begin{center}
\begin{minipage}{0.88\textwidth}
\kursiv{Die dritte Zeile ist eine harte Eigenschaft: Ein \code{panic} in irgendeiner Goroutine beendet das gesamte Programm --- es gibt keine Kapselung je Goroutine. Wer tausend Geräteverbindungen bedient, verliert alle, wenn eine davon einen Programmfehler auslöst.}
\end{minipage}
\end{center}

\begin{center}
\kursiv{Die Abhilfe ist ein Fangblock am Anfang jeder langlebigen Goroutine: \code{defer func() \{ if r := recover(); r != nil \{ melde(r) \} \}()}. Das ist eine der wenigen Stellen, an denen \code{recover} richtig ist.}
\end{center}
\end{warnung}
\end{merksatz}

\begin{aufgabe}
\begin{enumerate}[label=(\alph*)]
\item Rechnen Sie die Bauformen gegeneinander.
\item Begrenzen Sie die Zahl der Goroutinen.
\item Untersuchen Sie das Fadenmodell darunter.
\end{enumerate}

\loesung
\textbf{(a) Bauformen gegeneinander.}

\textbf{Die Aufgabe.} $1000$ Geräte, je $500$ Messwerte je Sekunde, Weiterleitung über das Netz.

\begin{center}
{\small
\renewcommand{\arraystretch}{1.45}
\begin{tabular}{@{}P{2.8cm}P{2.0cm}P{2.0cm}P{2.0cm}@{}}
\toprule
\textbf{Bauform} & \textbf{Zeilen} & \textbf{Speicher} & \textbf{Durchsatz}\\
\midrule
Ereignisschleife & $\mathbf{680}$ & \textcolor{gruen}{$14$ MB} & $\mathbf{510}$ k/s\\
Fadenbündel & $420$ & \textcolor{akzent}{$\approx400$ MB} & $180$ k/s\\
\textbf{Goroutinen} & $\mathbf{310}$ & $\mathbf{22}$ MB & $\mathbf{495}$ \textbf{k/s}\\
\bottomrule
\end{tabular}}
\end{center}

\begin{center}
\begin{minipage}{0.88\textwidth}
\kursiv{Die Goroutinen erreichen fast den Durchsatz der Ereignisschleife bei der Hälfte der Zeilen. Das ist der eigentliche Gewinn: Man schreibt den Ablauf hintereinander und bekommt die Leistung der zerrissenen Bauform.}
\end{minipage}
\end{center}

\textbf{Der Vergleich beim Lesen.}

\begin{center}
\renewcommand{\arraystretch}{1.45}
\begin{tabular}{@{}P{3.0cm}P{2.2cm}P{2.8cm}@{}}
\toprule
\textbf{Frage} & \textbf{Ereignis\-schleife} & \textbf{Goroutinen}\\
\midrule
Wo geht es weiter? & \textcolor{akzent}{\textbf{Rückruf suchen}} & \textcolor{gruen}{\textbf{nächste Zeile}}\\
Zustand je Gerät & \textcolor{akzent}{eigene Struktur} & \textcolor{gruen}{lokale Variablen}\\
Verstehenszeit & $\mathbf{175}$ s & \textcolor{gruen}{$\mathbf{48}$ s}\\
\bottomrule
\end{tabular}
\end{center}

\begin{center}
\kursiv{Faktor dreieinhalb bei der Verstehenszeit --- gemessen wie in Band 20, Kapitel~8. Der Grund steht in der zweiten Zeile: Bei Goroutinen ist der Zustand einer Verbindung in lokalen Variablen, nicht in einer von Hand geführten Struktur.}
\end{center}

\medskip
\textbf{(b) Die Zahl begrenzen.}

\textbf{Das Problem.} Eine Goroutine je eingehender Nachricht --- bei einer Lastspitze Millionen.

\begin{center}
\begin{minipage}{0.9\textwidth}
\code{sem := make(chan struct\{\}, 64)\ \ \ // hoechstens 64 zugleich}\\
\\
\code{for m := range eingang \{}\\
\code{\ \ \ \ sem <- struct\{\}\{\}}\\
\code{\ \ \ \ go func(m Messwert) \{}\\
\code{\ \ \ \ \ \ \ \ defer func() \{ <-sem \}()}\\
\code{\ \ \ \ \ \ \ \ verarbeite(m)}\\
\code{\ \ \ \ \}(m)}\\
\code{\}}
\end{minipage}
\end{center}

\begin{center}
\renewcommand{\arraystretch}{1.45}
\begin{tabular}{@{}P{2.4cm}P{2.2cm}P{2.4cm}P{2.0cm}@{}}
\toprule
\textbf{Grenze} & \textbf{Speicher} & \textbf{Durchsatz} & \textbf{Wartezeit}\\
\midrule
Keine & \textcolor{akzent}{\textbf{unbegrenzt}} & \textcolor{akzent}{bricht ein} & \textcolor{akzent}{$\infty$}\\
$8$ & \textcolor{gruen}{gering} & $180$ k/s & $44$ ms\\
$\mathbf{64}$ & \textcolor{gruen}{$\approx3$ MB} & $\mathbf{495}$ \textbf{k/s} & $\mathbf{13}$ \textbf{ms}\\
$512$ & $24$ MB & $501$ k/s & $19$ ms\\
\bottomrule
\end{tabular}
\end{center}

\begin{center}
\begin{minipage}{0.88\textwidth}
\kursiv{Die erste Zeile ist der Fehler, den Goroutinen so leicht machen: Weil sie billig sind, startet man beliebig viele --- und bei einer Lastspitze belegen sie mehr Speicher, als da ist. Die Begrenzung kostet fünf Zeilen und macht das System belastbar.}
\end{minipage}
\end{center}

\begin{center}
\kursiv{Bemerkenswert ist die letzte Zeile: Mehr Nebenläufigkeit bringt keinen Durchsatz mehr und erhöht die Wartezeit. Das Optimum liegt dort, wo alle Kerne beschäftigt sind und die Warteschlangen kurz bleiben --- gemessen, nicht geraten.}
\end{center}

\medskip
\textbf{(c) Das Fadenmodell darunter.}

\begin{center}
\renewcommand{\arraystretch}{1.5}
\begin{tabular}{@{}P{2.6cm}P{6.2cm}@{}}
\toprule
\textbf{Goroutinen} & Viele, von der Laufzeitumgebung verwaltet\\
\textbf{Systemfäden} & Wenige --- etwa so viele wie Kerne\\
\textbf{Die Zuordnung} & \textbf{Die Laufzeitumgebung verteilt Goroutinen auf Fäden}\\
\bottomrule
\end{tabular}
\end{center}

\begin{center}
\renewcommand{\arraystretch}{1.45}
\begin{tabular}{@{}P{3.4cm}P{2.2cm}P{2.6cm}@{}}
\toprule
\textbf{Vorgang} & \textbf{Faden blockiert?} & \\
\midrule
Kanaloperation & \textcolor{gruen}{\textbf{nein}} & \textbf{Goroutine wechselt}\\
Netzarbeit & \textcolor{gruen}{nein} & \textcolor{gruen}{eigener Melder}\\
Dateizugriff & \textcolor{mittel}{teils} & \\
\textbf{Aufruf nach C} & \textcolor{akzent}{\textbf{ja}} & \textcolor{akzent}{\textbf{Band 21, Kap.~14}}\\
Enge Rechenschleife & \textcolor{mittel}{unter\-brechbar} & seit Go 1.14\\
\bottomrule
\end{tabular}
\end{center}

\begin{center}
\begin{minipage}{0.88\textwidth}
\kursiv{Die vierte Zeile ist die Warnung aus Band 21: Ein blockierender C-Aufruf hält den Systemfaden an. Bei wenigen Aufrufen erzeugt die Laufzeitumgebung einen neuen Faden; bei vielen bricht die Nebenläufigkeit zusammen.}
\end{minipage}
\end{center}

\begin{center}
\kursiv{Die letzte Zeile war lange ein echtes Problem: Eine Rechenschleife ohne Funktionsaufruf konnte nicht unterbrochen werden und blockierte einen Kern. Seit Go 1.14 unterbricht die Laufzeitumgebung von sich aus --- eine der wenigen stillen Verbesserungen, die man bemerkt.}
\end{center}

\begin{warnung}
\boxtitelr{Billig ist nicht kostenlos}
\begin{center}
\renewcommand{\arraystretch}{1.5}
\begin{tabular}{@{}P{3.0cm}P{2.2cm}P{2.8cm}@{}}
\toprule
\textbf{Zahl} & \textbf{Speicher} & \textbf{Bewertung}\\
\midrule
$10^3$ & $\approx4$ MB & \textcolor{gruen}{unbedenklich}\\
$10^5$ & $\approx400$ MB & \textcolor{mittel}{spürbar}\\
$10^6$ & \textcolor{akzent}{$\approx\mathbf{4}$ \textbf{GB}} & \textcolor{akzent}{\textbf{Grenze}}\\
Nie endend & \textcolor{akzent}{wächst} & \textcolor{akzent}{\textbf{Leck}}\\
\bottomrule
\end{tabular}
\end{center}

\begin{center}
\begin{minipage}{0.88\textwidth}
\kursiv{Die letzte Zeile ist die gefährlichste und Gegenstand von Kapitel~\ref{kap:schiefgeht}: Eine Goroutine, die auf einen Kanal wartet, den niemand mehr bedient, bleibt für immer. Sie verbraucht Speicher und erscheint in keiner Statistik über Speicherlecks.}
\end{minipage}
\end{center}
\end{warnung}
\end{aufgabe}
\vspace{-4pt}
\noindent Pike\zit{18} trennt Nebenläufigkeit von Parallelität in einem Vortrag, der die Verwechslung beendet; Hoare\zit{13} legte 1978 die theo\-re\-ti\-sche Grundlage; die Vor\-trä\-ge zu Ne\-ben\-läu\-fig\-keits\-mus\-tern\zit{21} zeigen die wiederkehrenden Bauformen.
